% !TeX root = ../main.tex
% !TeX program = xelatex
% !TeX encoding = UTF-8
\section{结论}

本文以统一热量计量、连续建筑状态和物理约束连接供热规划与运行，并用分组实验解释不同决策的作用。

（1）固定代表日后，守恒约束与分布目标仍能改善压缩质量。六楼分布误差中位数由12.32\%降至4.48\%，统计检查通过数由3增至4。三城固定策略库的费用误差由0.5569\%降至0.1405\%，预算合格选择由9/12增至12/12。

（2）完整时序回放与机制对照承担不同任务。四楼历史数学场景完成8784小时服务年和2928小时声明冬季检查；厂家性能域下的三城六组调度降低电量费约1.10\%--1.60\%。双楼中，一楼热容增至四倍时，有、无区域上限的节费金额相差150.52元，表明单楼蓄热价值受共享容量影响。

（3）温压检查能够改变方案。参考管网三轮反馈将最大压降梯度降至237.96~Pa/m，最低远端水温保持57.45~$^\circ$C；相同交付热量下的费用变化可由热损与泵耗解释。

（4）三城中点碳帽保留最低电费方案80.92\%--89.44\%的节费额，并将其额外排放约减半。另在无碳帽、固定性能及$F=5$重复窗口条件下，单台采购预算切换点约9346.48元。两类分析分别为运行取舍和条件采购提供依据。

程序保存输入、配置和逐时结果，方案不仅能比较费用，还能追溯出力、室温及容量占用；替换具体项目参数后可沿用同一检查过程。
